% Generated draft: requires mathematical and source review.
\documentclass[11pt]{article}
\usepackage[margin=25mm]{geometry}
\usepackage{fontspec}
\setmainfont{Latin Modern Roman}
\usepackage{amsmath,amssymb,mathtools,unicode-math}
\setmathfont{Latin Modern Math}
\usepackage{enumitem,needspace,xurl,hyperref}
\setlength{\parindent}{0pt}
\setlength{\parskip}{5pt}
\setlength{\emergencystretch}{3em}
\newcommand{\N}{\mathbb{N}}
\newcommand{\Z}{\mathbb{Z}}
\newcommand{\Q}{\mathbb{Q}}
\newcommand{\R}{\mathbb{R}}
\newcommand{\C}{\mathbb{C}}
\newcommand{\F}{\mathbb{F}}
\newcommand{\A}{\mathbb{A}}
\newcommand{\PP}{\mathbb{P}}
\newcommand{\E}{\mathbb{E}}
\newcommand{\eps}{\varepsilon}
\newcommand{\dd}{\,\mathrm{d}}
\begin{document}
\section*{UnsolvedMath ID: 10400149}
\textbf{Source problem code:} catalogue entry
\subsection{Definitions and mathematical statement}
In the usual three-dimensional setting, a topological quantum field theory \((V,Z)\) assigns a finite-dimensional vector space \(V(\Sigma)\) to each closed surface and a linear map to each cobordism, compatibly with disjoint union, identities and gluing. Fix the intended coefficient field and cobordism category. Let \(\Sigma_g\) be a closed connected surface of genus \(g\).

The adjacent local source record supplies the definition omitted from this item:
\[
 P_{(V,Z)}(t)=\sum_{g=0}^{\infty}\dim V(\Sigma_g)\,t^g.
\]
It is a formal dimension series, with nonnegative integer coefficients in the ordinary finite-dimensional convention. Problem 8.3 asks:
\begin{enumerate}
\item Characterize exactly the formal series occurring as \(P_{(V,Z)}\) for the intended TQFTs.
\item For every series \(P(t)\) meeting that characterization, classify all TQFTs \((V,Z)\) with \(P_{(V,Z)}(t)=P(t)\), up to the intended theory equivalence.
\end{enumerate}
In the standard functorial convention the equivalence is a monoidal natural isomorphism. The second task concerns the full surface spaces and cobordism maps, not just their dimensions, so equal dimension series are not assumed to imply equivalent theories. The first task requires realizability of all coefficients by one theory, not separate genus-by-genus realizations.

The source does not fully specify the field, oriented or enhanced cobordisms, anomaly conventions or allowed TQFT class. The neighboring modular-category classification question does not itself restrict this task to semisimple, unitary, or modular-category-derived theories. No rational-function or Verlinde-form constraint is assumed without such hypotheses.
\subsection{Short English statement}
Characterize the genus-dimension series of three-dimensional TQFTs, then classify every theory realizing each allowed series. The definition is recovered from adjacent local context; the precise TQFT category and equivalence still need confirmation.
\subsection{Sources}
\begin{itemize}
\item Local UnsolvedMath registry, permanent ID 10400149, source code catalogue entry; no proposer or proposal year supplied; updated 2026-08-06.
\item Ohtsuki, \emph{Problems on invariants of knots and 3-manifolds} (2002), Problem 8.3, PDF page 123: \url{https://msp.org/gtm/2002/04/gtm-2002-04-024s.pdf}. Registry access date 2026-07-29; PDF font-structure extraction. Provenance distinguishes 231 numbered bold headings. Present primary-literature status and redistribution terms require review.
\item Adjacent local ID 10400148, source code AMR-103-0148, Problem 8.2, PDF page 122, explicitly defines the dimension series and calls the next problem a refinement of Problem 8.1. No external source was visited. The supplied summary adds no precise resolving reference.
\end{itemize}
\end{document}
